\makereporttitle{2024}{iulie}

\makereportpreamble

\textbf{Activitate de cercetare:}

În această lună am parcurs diferite articole de fizică și de fizică matematică, în special pe temele discutate la Școala Regională de Matematici Aplicate (menționată mai jos).

\textbf{Cursuri și seminare științifice:}

Între 10 și 18 iulie am participat la a 9-a ediție a \href{https://sites.google.com/site/marinliviu/regional-french-romanian-summer-school-on-applied-mathematics/9th-regional-summer-school-on-applied-mathematics-10-18-july-2024?authuser=0}{Regional Summer School on Applied Mathematics}, care s-a desfășurat la Stațiunea Zoologică Sinaia a Universității din București. La această școală au fost susținute prelegeri pe diferite teme științifice:

\begin{itemize}
    \item \emph{Recent Advances on Continuous Optimization and Applications}, de către Ion Necoară (asociat cu Universitatea Națională de Știință și Tehnologie POLITEHNICA București și cu Institutul de Statistică Matematică și Matematică Aplicată „Gheorghe Mihoc - Caius Iacob” al Academiei Române);

    \item \emph{Old and New in Homogenisation, with Applications to Liquid Crystals}, de către Arghir Zărnescu (Centrul Basc pentru Matematică Aplicată și Institutul de Matematică „Simion Stoilow” al Academiei Române);

    \item \emph{Wave Propagation in Heterogeneous Media}, de către Ioan R. Ionescu (Universitatea Sorbonne --- Paris-Nord);

    \item \emph{A Panorama of Rough Analysis}, de către Nicolas Perowski (Universitatea Liberă din Berlin);
\end{itemize}

De asemenea, am participat și la o conferință științifică organizată de FMI-IMAR, intitulată \emph{Stochastic gradient descent optimization methods with adaptive learning rates}, susținută de prof.~dr.~Arnulf Jentzen.

\vspace{0.5em}

\textbf{Observațiile conducătorului științific}:

\vspace{3em}

\makereportpostamble{29 iulie 2024}
